% ============================================================
%  From Reports to Decisions: Evaluating LLM-Based Traffic
%  Incident Summarization for Smart City Operations
% ============================================================

\documentclass[10pt,conference]{IEEEtran}

% ── Encoding & fonts ─────────────────────────────────────────
\usepackage[T1]{fontenc}
\usepackage[utf8]{inputenc}

% ── Core packages ────────────────────────────────────────────
\usepackage{cite}
\usepackage{amsmath,amssymb}
\usepackage{booktabs}
\usepackage{array}
\usepackage[table]{xcolor}
\usepackage{url}
% hyperref removed for IEEE PDF eXpress compliance (no bookmarks, no links)
\newcommand{\href}[2]{#2}

% ── Table helpers ────────────────────────────────────────────
\usepackage{multirow}
\usepackage{tabularx}
\usepackage{makecell}

% ── Layout tuning ────────────────────────────────────────────
\tolerance=1000
\hbadness=10000
\def\IEEEkeywordsname{Keywords}
% Compact section spacing
\usepackage[compact]{titlesec}

% ── Custom column type ───────────────────────────────────────
\newcolumntype{C}[1]{>{\centering\arraybackslash}p{#1}}

% ── Table colors ────────────────────────────────────────────
\definecolor{HeaderBlue}{RGB}{184,204,228}
\definecolor{SectionBlue}{RGB}{217,226,243}
\definecolor{RowYellow}{RGB}{242,242,230}
\definecolor{RowGray}{RGB}{242,242,242}

% ============================================================
\begin{document}

\title{Automated Traffic Incident Summarisation Using LLMs: A Smart City Perspective}

\author{
\IEEEauthorblockN{Rajeev Ranjan Pandey\IEEEauthorrefmark{1}, Khalid Elgazzar\IEEEauthorrefmark{2}\IEEEauthorrefmark{1}}
\IEEEauthorblockA{\IEEEauthorrefmark{1}20250007918@students.cud.ac.ae, \IEEEauthorrefmark{2}khalid.elgazzar@ontariotechu.ca}
\IEEEauthorblockA{
\IEEEauthorrefmark{1}Canadian University Dubai, UAE\\
\IEEEauthorrefmark{2}IoT Research Laboratory, Ontario Tech University, Canada}
}

\maketitle

% ============================================================
\begin{abstract}
Traffic control centers process incident reports that are often too long for rapid operational use. This paper evaluates transformer-based summarization for converting such reports into concise, decision-ready briefings. Using a dataset of over 5,250 records derived from the US Accidents dataset and structured Dubai Pulse traffic data converted into narrative descriptions, BART-large-cnn, Flan-T5-small, and PEGASUS-CNN/DailyMail are compared with TextRank under zero-shot settings. BART achieves the strongest overall performance with ROUGE-1 of 0.432, ROUGE-2 of 0.198, and ROUGE-L of 0.391, outperforming TextRank by 35.8\% on ROUGE-1. However, qualitative analysis shows that strong lexical scores do not ensure operational reliability: BART hallucinates safety-critical details in approximately 14\% of outputs, while Flan-T5 omits road-closure information in approximately 11\%. These findings show that traffic incident summarization systems for smart city operations should be evaluated using both lexical overlap and factual consistency.
\end{abstract}

\begin{IEEEkeywords}
traffic incident summarization, abstractive summarization, transformer models, zero-shot evaluation, ROUGE metrics, factual hallucination, smart city operations, GCC traffic data, safety-critical NLP, extractive-abstractive comparison
\end{IEEEkeywords}

% ============================================================

\section{Introduction}
\vspace{-1.5mm}

Smart cities generate continuous streams of traffic incident reports for
control rooms, dispatch, and mobility monitoring. These reports are too long for fast decision-making, yet dispatchers must scan lengthy narratives to find the incident type, road segment, closure status, and severity. Automatic summarization can bridge this gap. Extractive methods are simple but miss distributed facts; abstractive transformer models generate coherent summaries but may hallucinate unsafe details.

\textbf{Contributions:}
(1) A dataset of over 5,250 records combining US Accidents~\cite{moosavi2023} and Dubai Pulse~\cite{dubaipulse2023}, providing GCC regional grounding.
(2) A zero-shot evaluation of BART, Flan-T5, and PEGASUS against TextRank using ROUGE and qualitative failure analysis.
(3) A live comparison tool for applied smart mobility contexts.

% ============================================================
\section{Related Work}
\vspace{-1.5mm}

BART~\cite{lewis2020} uses denoising seq2seq pretraining for strong contextual modeling. T5~\cite{raffel2020} formulates all NLP tasks as text-to-text, making instruction-style summarization natural. PEGASUS~\cite{zhang2020} uses gap-sentence generation, a pretraining objective designed for summarization. TextRank is the extractive baseline, ranking sentences by graph centrality; it preserves factual wording and avoids hallucination but misses distributed facts. A key limitation of abstractive models is factual reliability: Maynez et al.\ showed that high ROUGE scores do not guarantee faithfulness. This is especially critical in traffic operations where wrong closure or injury information affects dispatch decisions.

\textbf{Proposed Approach.}
This study adopts a zero-shot framework evaluating BART, Flan-T5, and PEGASUS against TextRank. In addition to ROUGE, qualitative failure analysis captures hallucination, omission, and over-compression errors, making the evaluation suitable for safety-critical deployments.

% ============================================================
\section{Methodology}
\vspace{-1mm}
\subsection{Dataset Construction}
\vspace{-1mm}

The corpus combines over 5,250 records from the US Accidents
dataset~\cite{moosavi2023} with GCC data from Dubai Pulse, UAE Federal
Traffic Statistics, and Abu Dhabi Open Data~\cite{dubaipulse2023}.
Because GCC portals mainly provide structured records, these were
normalized and converted into operator-style narrative incident reports
using a rule-based text generation approach. This resulted in 250 GCC narrative samples aligned
with operational reporting formats.

ROUGE evaluation used a stratified set of 420 records (300 US, 120 GCC)
across five incident categories: rear-end collisions, highway lane
blockages, weather-related events, pedestrian incidents, and
multi-vehicle pile-ups. Each record was paired with a human-written
reference summary of 40 to 80 tokens, allowing both quantitative and
qualitative comparison across models.

\subsection{Models and Configuration}
\vspace{-1mm}

BART-large-cnn~\cite{lewis2020}, Flan-T5-small~\cite{raffel2020}, and
PEGASUS-CNN/DailyMail~\cite{zhang2020} were evaluated in zero-shot mode.
BART was selected because incident reports share some structural
similarity with news-style event descriptions. Flan-T5 was included as a
lighter instruction-oriented model that supports prompt-based
summarization through the \texttt{summarize:} prefix. PEGASUS was used as
a summarization-specialized model trained with a gap-sentence objective,
providing an additional abstractive comparison point. TextRank served as
the extractive baseline.

Inputs were truncated to 512 tokens; beam size 2, length penalty 1.0, max-new-tokens 72, min-new-tokens 18, no-repeat-ngram-size 3. ROUGE metrics were computed using the rouge-score library with stemming. Execution times were recorded on a standard cloud GPU environment.

% ============================================================
\section{Experimental Results and Analysis}
\vspace{-1mm}

Table~\ref{tab:rouge} presents ROUGE scores, compression ratios (CR),
execution times, and ROUGE-1 gains over TextRank. All three abstractive models outperform
the baseline. BART achieves the strongest overall results
(ROUGE-1: 0.432, ROUGE-2: 0.198, ROUGE-L: 0.391), representing a 35.8\%
improvement over TextRank.

\begin{table}[htbp]
\caption{ROUGE Scores (R-1/R-2/R-L), Compression Ratio (CR),
ROUGE-1 Gain vs.\ TextRank ($n = 420$), and Execution Time}
\label{tab:rouge}
\centering
\renewcommand{\arraystretch}{1.1}
\setlength{\tabcolsep}{3pt}
\arrayrulecolor{gray}
\begin{tabular}{|p{2.25cm}|C{0.6cm}|C{0.6cm}|C{0.6cm}|C{0.55cm}|C{0.95cm}|C{1.1cm}|}
\hline
\rowcolor{HeaderBlue}
\rule{0pt}{2.4ex}\textbf{Model} & \textbf{R-1} & \textbf{R-2} & \textbf{R-L} & \textbf{CR} & \textbf{Gain R-1} & \textbf{Exec.(s)} \\
\hline
\rowcolor{SectionBlue}
\multicolumn{7}{|l|}{\textit{\textbf{Abstractive Models}}}\\
\hline
\rowcolor{RowYellow}
\textbf{BART-large-cnn} & \textbf{0.432} & \textbf{0.198} & \textbf{0.391} & 3.2x & \textbf{+35.8\%} & 1.42 \\
\hline
Flan-T5-small & 0.408 & 0.181 & 0.372 & 3.5x & +28.3\% & 0.28 \\
\hline
PEGASUS-CNN/DM & 0.389 & 0.162 & 0.354 & 3.8x & +22.3\% & 1.65 \\
\hline
\rowcolor{SectionBlue}
\multicolumn{7}{|l|}{\textit{\textbf{Extractive Baseline}}}\\
\hline
\rowcolor{RowGray}
TextRank & 0.318 & 0.109 & 0.287 & 2.8x & --- & \textbf{0.05} \\
\hline
\end{tabular}
\\[2pt]
{\small\textit{R-1 = ROUGE-1, R-2 = ROUGE-2, R-L = ROUGE-L, CR = Compression Ratio, Exec.(s) = Exec. Time per record.}}
\end{table}

All abstractive models outperform TextRank. BART performs best on lexical metrics, while Flan-T5 offers a faster (0.28 s/record) trade-off for time-sensitive deployments versus BART's 1.42 s/record. PEGASUS produces over-compressed outputs less suited to multi-fact traffic reporting.

\textbf{Error Analysis.}
BART hallucinates details in $\approx$14\% of outputs (e.g., fabricating injury counts); mitigation includes constraint-based decoding or factuality checks. Flan-T5 omits road closure information in $\approx$11\% of cases. PEGASUS over-compresses, removing location or response details. TextRank avoids hallucination but misses facts distributed across multiple sentences.

ROUGE alone is insufficient for safety-critical systems. Future evaluations should incorporate BERTScore or similar semantic metrics, and hybrid extractive-abstractive designs to balance factual grounding and fluency.

% ============================================================
\section{Conclusion}
\vspace{-1mm}

This study evaluates transformer-based summarization for traffic incident reports. BART achieves the best ROUGE performance but hallucinates in $\approx$14\% of outputs. Future systems should combine extractive grounding with abstractive generation, domain-specific fine-tuning, Arabic-language support for GCC, and semantic evaluation via BERTScore.

\vspace{-1mm}
\section*{Code \& Demo}
Code: GitHub (VizArtPandey/CUD\_AAI\_MINIPROJECT\_LLMs-for-Traffic-Incident-Summarization). Demo: Hugging Face Spaces (rajvivan/CUD-Traffic-AI).

% ============================================================
\vspace{-2mm}
\begin{thebibliography}{7}
\footnotesize

\bibitem{lewis2020}
M.~Lewis, Y.~Liu, N.~Goyal, M.~Ghazvininejad, A.~Mohamed, O.~Levy,
V.~Stoyanov, and L.~Zettlemoyer,
``BART: Denoising Sequence-to-Sequence Pre-training for Natural Language
Generation, Translation, and Comprehension,''
in \textit{Proc.\ 58th Annu.\ Meeting ACL}, Online, 2020, pp.\ 7871--7880.

\bibitem{raffel2020}
C.~Raffel, N.~Shazeer, A.~Roberts, K.~Lee, S.~Narang, M.~Matena,
Y.~Zhou, W.~Li, and P.~J.~Liu,
``Exploring the Limits of Transfer Learning with a Unified
Text-to-Text Transformer,''
\textit{J.\ Mach.\ Learn.\ Res.}, vol.~21, no.~140, pp.~1--67, 2020.

\bibitem{zhang2020}
J.~Zhang, Y.~Zhao, M.~Saleh, and P.~J.~Liu,
``PEGASUS: Pre-training with Extracted Gap-sentences for Abstractive
Summarization,''
in \textit{Proc.\ ICML}, Vienna, 2020, pp.~11328--11339.

\bibitem{moosavi2023}
S.~Moosavi Nezhad Asl,
``US Accidents: A Countrywide Traffic Accident Dataset (2016--2023),''
Kaggle, 2023. [Online]. Available:
\url{https://www.kaggle.com/datasets/sobhanmoosavi/us-accidents}



\bibitem{dubaipulse2023}
Dubai Statistics Center,
``Dubai Pulse Open Data Portal,''
Government of Dubai, 2023. [Online]. Available:
\url{https://www.dubaipulse.gov.ae}

\end{thebibliography}

\end{document}